\section{模型假设}

为简化问题，本文做出以下假设：

假设1：典型日的电价、小区负载与光伏发电功率在每个10分钟时段内保持不变，且逐日完全重复，即问题一中的优化是在完全信息下进行的确定性决策。

假设2：将题目中的“充放电效率为90\%”按充、放电两个环节的单程效率理解，即储能设备的充电效率与放电效率均恒为0.9且不随储电量与功率变化；充放电过程不存在自放电损耗，储电量在0:00至24:00之间始终处于1200至10800 kWh的允许区间内。

假设3：光伏发电与小区负载的预测误差具有短期平稳性，可用最近若干个完整日的实测误差样本刻画其分布，调度中心仅使用决策时已完成时段的观测与已经发布的预报。

假设4：微网与外网之间的功率交换不受联络线容量限制，计划与调整仅在题定时刻提交，剩余供电缺口允许紧急采购；不向外网售电。问题四的未来价格视为未知，决策用历史价格预测，实际结算用附件4实现价。
